% chapter10_modern_entropy_coders.tex
\section{Modern Entropy Coders — ANS and Range Coding}

\subsection{Introduction: Beyond Huffman's Limitations}

Entropy coding is the final stage of most modern compression systems.  
Earlier chapters introduced two fundamental approaches:

\begin{itemize}
    \item \textbf{Huffman Coding} (Chapter 2.8) — fast but limited to integer-bit codeword lengths.
    \item \textbf{Arithmetic Coding} (Chapter 3.7) — achieves fractional-bit precision and approaches the Shannon limit.
\end{itemize}

Although arithmetic coding provides near-optimal compression, it has historically faced several challenges:

\begin{enumerate}
    \item \textbf{Patent restrictions} (expired around 2005).
    \item \textbf{High computational cost}, especially multiplications and divisions.
    \item \textbf{Complex implementation}, particularly renormalization and interval management.
\end{enumerate}

These limitations motivated the search for faster, simpler entropy coders that still preserve the compression efficiency of arithmetic coding.  
This chapter presents two such techniques:

\begin{itemize}
    \item \textbf{Range Coding} — a practical integer formulation of arithmetic coding used in LZMA, WebP, and ffmpeg.
    \item \textbf{Asymmetric Numeral Systems (ANS)} — a breakthrough method discovered in 2009 that combines Huffman-like speed with arithmetic-like compression.
\end{itemize}

Both techniques are fundamental to modern compressors, including Zstd, LZFSE, JPEG~XL, and other industry standards.

\subsection{Range Coding: Arithmetic With Integers}

Range coding is best understood as a re-interpretation of arithmetic coding where the ``interval'' is represented using integers rather than fractions.  
This gives several advantages:

\begin{itemize}
    \item simpler renormalization,
    \item integer-only arithmetic,
    \item complete avoidance of arithmetic coding patents,
    \item highly efficient implementations (used by LZMA, WebP, VPx, and others).
\end{itemize}

\subsubsection{Core Concept}

Arithmetic coding encodes a message by narrowing a fractional interval $[0,1)$.  
Range coding reformulates this idea:

\begin{quote}
\emph{Represent the current interval as integer values \texttt{low} and \texttt{range} where the actual interval is $[\,\texttt{low},\, \texttt{low} + \texttt{range}\,)$ inside a fixed-size integer space (usually 32 bits).}
\end{quote}

For each symbol $s$ with frequency $f(s)$ and cumulative frequency $F(s)$ (where $\sum f(s) = \text{total}$), the interval is updated as:

\begin{algorithm}[H]
\caption{Modern Range Coding Update}
\begin{algorithmic}[1]
\Procedure{EncodeSymbol}{$s$, $freq$, $cum$, $total$}
    \State $\text{range} \gets \text{range} \times \text{freq}[s] / \text{total}$  \Comment{First narrow range}
    \State $\text{low} \gets \text{low} + \text{range} \times \text{cum}[s] / \text{total}$  \Comment{Then adjust low}
    \State \Call{Renormalize}{}
\EndProcedure
\end{algorithmic}
\end{algorithm}

Most modern implementations use \textbf{scaled integer frequencies}, where \texttt{total} is a power of two (e.g., $2^{12}$, $2^{15}$) to replace division with bit shifts.

\subsubsection{Renormalization with Byte Output}

Because \texttt{low} and \texttt{range} are stored in a fixed-width integer (e.g., 32 bits), we must periodically \emph{renormalize} to maintain precision. The industry standard (LZMA, WebP) uses:

\[
\text{threshold} = 2^{24} = 16{,}777{,}216
\]

Whenever $\text{range} < 2^{24}$, we output the most significant byte of \texttt{low} and shift:

\begin{algorithm}[H]
\caption{Renormalization (Industry Standard)}
\begin{algorithmic}[1]
\Procedure{Renormalize}{}
    \While{$\text{range} < 2^{24}$}
        \State $\text{output\_byte}(\text{low} \gg 24)$
        \State $\text{low} \leftarrow (\text{low} \ll 8) \;\&\; \text{0xFFFFFFFF}$
        \State $\text{range} \leftarrow \text{range} \ll 8$
    \EndWhile
\EndProcedure
\end{algorithmic}
\end{algorithm}

This ensures at least 24 bits of precision at all times and produces a byte-aligned output stream.

\subsubsection{Advantages}

\begin{itemize}
    \item No floating-point operations.
    \item The method is patent-free.
    \item Excellent performance on modern CPUs.
    \item Simple to integrate into block compressors.
\end{itemize}

\subsubsection{Range Coding Example with Renormalization}

Consider an alphabet $\{A,B,C\}$ with frequencies:

\[
f(A) = 5,\quad f(B) = 3,\quad f(C) = 2,\quad \text{total} = 10.
\]

Cumulative frequencies:

\[
c(A) = 0,\quad c(B) = 5,\quad c(C) = 8.
\]

Symbol intervals within $[0,10)$:

\[
A: [0,5),\quad B: [5,8),\quad C: [8,10).
\]

We encode the sequence: \texttt{B A C C C C} (six symbols total). This sequence includes different symbols and enough $C$'s (the smallest frequency) to trigger renormalization.

\paragraph{Encoding Algorithm}

Initialize with 32-bit precision:
\[
\text{low} = 0,\qquad \text{range} = 2^{32} = 4{,}294{,}967{,}296.
\]

For each symbol $x$ with cumulative frequency $c_x$ and frequency $f_x$:

\begin{enumerate}
    \item Update low: $\text{low} \gets \text{low} + \text{range} \times \dfrac{c_x}{\text{total}}$
    \item Update range: $\text{range} \gets \text{range} \times \dfrac{f_x}{\text{total}}$
    \item Renormalize: while $\text{range} < 2^{24} = 16{,}777{,}216$, output one byte and scale
\end{enumerate}

\emph{Important: The $\text{low}$ update uses the \textbf{old} range value.}

\paragraph{Step 1: Encode 'B'}

For 'B': $c_B = 5$, $f_B = 3$

\[
\begin{aligned}
\text{low} &\gets 0 + 4{,}294{,}967{,}296 \times \frac{5}{10} = 2{,}147{,}483{,}648 \\
\text{range} &\gets 4{,}294{,}967{,}296 \times \frac{3}{10} = 1{,}288{,}490{,}188
\end{aligned}
\]

Renormalization check: $1{,}288{,}490{,}188 \ge 16{,}777{,}216$, so continue.

\paragraph{Step 2: Encode 'A'}

For 'A': $c_A = 0$, $f_A = 5$

\[
\begin{aligned}
\text{low} &\gets 2{,}147{,}483{,}648 + 1{,}288{,}490{,}188 \times \frac{0}{10} = 2{,}147{,}483{,}648 \\
\text{range} &\gets 1{,}288{,}490{,}188 \times \frac{5}{10} = 644{,}245{,}094
\end{aligned}
\]

Renormalization check: $644{,}245{,}094 \ge 16{,}777{,}216$, so continue.

\paragraph{Step 3: Encode 'C'}

For 'C': $c_C = 8$, $f_C = 2$

\[
\begin{aligned}
\text{low} &\gets 2{,}147{,}483{,}648 + 644{,}245{,}094 \times \frac{8}{10} \\
&= 2{,}147{,}483{,}648 + 515{,}396{,}075 = 2{,}662{,}879{,}723 \\
\text{range} &\gets 644{,}245{,}094 \times \frac{2}{10} = 128{,}849{,}018
\end{aligned}
\]

Renormalization check: $128{,}849{,}018 \ge 16{,}777{,}216$, so continue.

\paragraph{Step 4: Encode 'C'}

For 'C': $c_C = 8$, $f_C = 2$

\[
\begin{aligned}
\text{low} &\gets 2{,}662{,}879{,}723 + 128{,}849{,}018 \times \frac{8}{10} \\
&= 2{,}662{,}879{,}723 + 103{,}079{,}214 = 2{,}765{,}958{,}937 \\
\text{range} &\gets 128{,}849{,}018 \times \frac{2}{10} = 25{,}769{,}803
\end{aligned}
\]

Renormalization check: $25{,}769{,}803 \ge 16{,}777{,}216$, so continue.

\paragraph{Step 5: Encode 'C'}

For 'C': $c_C = 8$, $f_C = 2$

\[
\begin{aligned}
\text{low} &\gets 2{,}765{,}958{,}937 + 25{,}769{,}803 \times \frac{8}{10} \\
&= 2{,}765{,}958{,}937 + 20{,}615{,}842 = 2{,}786{,}574{,}779 \\
\text{range} &\gets 25{,}769{,}803 \times \frac{2}{10} = 5{,}153{,}960
\end{aligned}
\]

Renormalization check: $5{,}153{,}960 < 16{,}777{,}216$ — \textbf{renormalization triggered!}

\paragraph{Renormalization (after Step 5)}

We output bytes until $\text{range} \ge 2^{24}$.

\textbf{Iteration 1:}
\[
\begin{aligned}
\text{output} &\gets \text{low} \gg 24 = 2{,}786{,}574{,}779 \gg 24 \\
2{,}786{,}574{,}779 &= 0\text{xA61A2FBB} \\
0\text{xA61A2FBB} \gg 24 &= 0\text{xA6} = 166 \\
\text{low} &\gets (\text{low} \ll 8) \mathrel{\&} (2^{32}-1) \\
&= (2{,}786{,}574{,}779 \ll 8) \mathrel{\&} 0\text{xFFFFFFFF} \\
&= 0\text{xA61A2FBB00} \mathrel{\&} 0\text{xFFFFFFFF} = 0\text{x1A2FBB00} = 439{,}463{,}680 \\
\text{range} &\gets \text{range} \ll 8 = 5{,}153{,}960 \times 256 = 1{,}319{,}413{,}760
\end{aligned}
\]

Now $\text{range} = 1{,}319{,}413{,}760 \ge 16{,}777{,}216$, so stop renormalization.

Output byte: $\text{0xA6}$ (166).

\paragraph{Step 6: Encode 'C'}

For 'C': $c_C = 8$, $f_C = 2$

\[
\begin{aligned}
\text{low} &\gets 439{,}463{,}680 + 1{,}319{,}413{,}760 \times \frac{8}{10} \\
&= 439{,}463{,}680 + 1{,}055{,}531{,}008 = 1{,}494{,}994{,}688 \\
\text{range} &\gets 1{,}319{,}413{,}760 \times \frac{2}{10} = 263{,}882{,}752
\end{aligned}
\]

Renormalization check: $263{,}882{,}752 \ge 16{,}777{,}216$, so continue.

\paragraph{Final Encoding Output}

The compressed representation consists of:
\begin{itemize}
    \item The byte $\text{0xA6}$ output during renormalization
    \item The final $\text{low} = 1{,}494{,}994{,}688$ (output as 4 bytes)
\end{itemize}

Thus the encoded stream is: $[166, 1{,}494{,}994{,}688]$ (5 bytes total).

In hexadecimal: $[\text{0xA6}, \text{0x59}, \text{0x1C}, \text{0x00}, \text{0x00}]$.

\paragraph{Encoding State Table}

\begin{center}
\begin{tabular}{c|c|c|c|c}
\hline
\textbf{Step} & \textbf{Symbol} & \textbf{Low} & \textbf{Range} & \textbf{Renorm Output} \\ \hline
Start & — & $0$ & $4{,}294{,}967{,}296$ & — \\
1 & B & $2{,}147{,}483{,}648$ & $1{,}288{,}490{,}188$ & — \\
2 & A & $2{,}147{,}483{,}648$ & $644{,}245{,}094$ & — \\
3 & C & $2{,}662{,}879{,}723$ & $128{,}849{,}018$ & — \\
4 & C & $2{,}765{,}958{,}937$ & $25{,}769{,}803$ & — \\
5 & C & $2{,}786{,}574{,}779$ & $5{,}153{,}960$ & $166$ \\
6 & C & $1{,}494{,}994{,}688$ & $263{,}882{,}752$ & — \\ \hline
\end{tabular}
\end{center}

\paragraph{Visualization of Interval Narrowing}

The table below shows how the interval $[\text{low}, \text{low}+\text{range})$ narrows with each symbol:

\begin{center}
\begin{tabular}{c|c|c|c}
\hline
\textbf{Step} & \textbf{Interval} & \textbf{Width} & \textbf{Symbol} \\ \hline
Start & $[0, 4.295\times10^9)$ & $4.295\times10^9$ & — \\
After B & $[2.147\times10^9, 3.436\times10^9)$ & $1.288\times10^9$ & B \\
After A & $[2.147\times10^9, 2.792\times10^9)$ & $6.442\times10^8$ & A \\
After C & $[2.663\times10^9, 2.792\times10^9)$ & $1.288\times10^8$ & C \\
After C & $[2.766\times10^9, 2.792\times10^9)$ & $2.577\times10^7$ & C \\
After C & $[2.787\times10^9, 2.792\times10^9)$ & $5.154\times10^6$ & C \\
Renorm & $[4.395\times10^8, 1.759\times10^9)$ & $1.319\times10^9$ & — \\
After C & $[1.495\times10^9, 1.759\times10^9)$ & $2.639\times10^8$ & C \\ \hline
\end{tabular}
\end{center}

\newpage

\paragraph{Decoding Algorithm}

Given the encoded stream $[166, 1{,}494{,}994{,}688]$ (bytes: $\text{0xA6}, \text{0x59}, \text{0x1C}, \text{0x00}, \text{0x00}$), we reconstruct the original sequence.

Initialize:
\[
\text{low} = 0,\qquad \text{range} = 2^{32} = 4{,}294{,}967{,}296
\]

Read the first 4 bytes from the stream as the initial code:
\[
\text{code} = 0\text{xA6591C00} = 2{,}789{,}021{,}696
\]

\paragraph{Step 1: Decode First Symbol}

Current interval: $[0, 4{,}294{,}967{,}296)$

Compute offset:
\[
u = \frac{\text{code} - \text{low}}{\text{range}} = \frac{2{,}789{,}021{,}696}{4{,}294{,}967{,}296} \approx 0.6494
\]

Find symbol containing $0.6494$:
\begin{itemize}
    \item $A$: $[0, 0.5)$ — no
    \item $B$: $[0.5, 0.8)$ — contains $0.6494$? Yes
    \item $C$: $[0.8, 1.0)$ — no
\end{itemize}
First symbol is $B$.

Update for $B$ ($c_B=5$, $f_B=3$):
\[
\begin{aligned}
\text{low} &\gets 0 + 4{,}294{,}967{,}296 \times \frac{5}{10} = 2{,}147{,}483{,}648 \\
\text{range} &\gets 4{,}294{,}967{,}296 \times \frac{3}{10} = 1{,}288{,}490{,}188
\end{aligned}
\]

Renormalization check: $\text{range} = 1{,}288{,}490{,}188 \ge 16{,}777{,}216$, so continue.

\paragraph{Step 2: Decode Second Symbol}

Current interval: $[2{,}147{,}483{,}648, 3{,}435{,}973{,}836)$

Compute offset:
\[
u = \frac{2{,}789{,}021{,}696 - 2{,}147{,}483{,}648}{1{,}288{,}490{,}188} = \frac{641{,}538{,}048}{1{,}288{,}490{,}188} \approx 0.498
\]

Find symbol containing $0.498$:
\begin{itemize}
    \item $A$: $[0, 0.5)$ — contains $0.498$? Yes
    \item $B$: $[0.5, 0.8)$ — no
    \item $C$: $[0.8, 1.0)$ — no
\end{itemize}
Second symbol is $A$.

Update for $A$ ($c_A=0$, $f_A=5$):
\[
\begin{aligned}
\text{low} &\gets 2{,}147{,}483{,}648 + 1{,}288{,}490{,}188 \times \frac{0}{10} = 2{,}147{,}483{,}648 \\
\text{range} &\gets 1{,}288{,}490{,}188 \times \frac{5}{10} = 644{,}245{,}094
\end{aligned}
\]

Renormalization check: $\text{range} = 644{,}245{,}094 \ge 16{,}777{,}216$, so continue.

\paragraph{Step 3: Decode Third Symbol}

Current interval: $[2{,}147{,}483{,}648, 2{,}791{,}728{,}742)$

Compute offset:
\[
u = \frac{2{,}789{,}021{,}696 - 2{,}147{,}483{,}648}{644{,}245{,}094} = \frac{641{,}538{,}048}{644{,}245{,}094} \approx 0.9958
\]

Find symbol containing $0.9958$:
\begin{itemize}
    \item $A$: $[0, 0.5)$ — no
    \item $B$: $[0.5, 0.8)$ — no
    \item $C$: $[0.8, 1.0)$ — contains $0.9958$? Yes
\end{itemize}
Third symbol is $C$.

Update for $C$ ($c_C=8$, $f_C=2$):
\[
\begin{aligned}
\text{low} &\gets 2{,}147{,}483{,}648 + 644{,}245{,}094 \times \frac{8}{10} \\
&= 2{,}147{,}483{,}648 + 515{,}396{,}075 = 2{,}662{,}879{,}723 \\
\text{range} &\gets 644{,}245{,}094 \times \frac{2}{10} = 128{,}849{,}018
\end{aligned}
\]

Renormalization check: $\text{range} = 128{,}849{,}018 \ge 16{,}777{,}216$, so continue.

\paragraph{Step 4: Decode Fourth Symbol}

Current interval: $[2{,}662{,}879{,}723, 2{,}791{,}728{,}741)$

Compute offset:
\[
u = \frac{2{,}789{,}021{,}696 - 2{,}662{,}879{,}723}{128{,}849{,}018} = \frac{126{,}141{,}973}{128{,}849{,}018} \approx 0.9790
\]

Find symbol containing $0.9790$:
\begin{itemize}
    \item $A$: $[0, 0.5)$ — no
    \item $B$: $[0.5, 0.8)$ — no
    \item $C$: $[0.8, 1.0)$ — contains $0.9790$? Yes
\end{itemize}
Fourth symbol is $C$.

Update for $C$ ($c_C=8$, $f_C=2$):
\[
\begin{aligned}
\text{low} &\gets 2{,}662{,}879{,}723 + 128{,}849{,}018 \times \frac{8}{10} \\
&= 2{,}662{,}879{,}723 + 103{,}079{,}214 = 2{,}765{,}958{,}937 \\
\text{range} &\gets 128{,}849{,}018 \times \frac{2}{10} = 25{,}769{,}803
\end{aligned}
\]

Renormalization check: $\text{range} = 25{,}769{,}803 \ge 16{,}777{,}216$, so continue.

\paragraph{Step 5: Decode Fifth Symbol}

Current interval: $[2{,}765{,}958{,}937, 2{,}791{,}728{,}740)$

Compute offset:
\[
u = \frac{2{,}789{,}021{,}696 - 2{,}765{,}958{,}937}{25{,}769{,}803} = \frac{23{,}062{,}759}{25{,}769{,}803} \approx 0.8950
\]

Find symbol containing $0.8950$:
\begin{itemize}
    \item $A$: $[0, 0.5)$ — no
    \item $B$: $[0.5, 0.8)$ — no
    \item $C$: $[0.8, 1.0)$ — contains $0.8950$? Yes
\end{itemize}
Fifth symbol is $C$.

Update for $C$ ($c_C=8$, $f_C=2$):
\[
\begin{aligned}
\text{low} &\gets 2{,}765{,}958{,}937 + 25{,}769{,}803 \times \frac{8}{10} \\
&= 2{,}765{,}958{,}937 + 20{,}615{,}842 = 2{,}786{,}574{,}779 \\
\text{range} &\gets 25{,}769{,}803 \times \frac{2}{10} = 5{,}153{,}960
\end{aligned}
\]

Now $\text{range} = 5{,}153{,}960 < 16{,}777{,}216$ — \textbf{renormalization triggered during decoding!}

\paragraph{Renormalization During Decoding}

When $\text{range} < 2^{24}$ during decoding, we:
\begin{enumerate}
    \item Update low: $\text{low} \gets (\text{low} \ll 8) \mathrel{\&} (2^{32}-1)$
    \item Read next byte from stream and incorporate into code: $\text{code} \gets ((\text{code} \ll 8) \mathrel{\&} (2^{32}-1)) \mid \text{next\_byte}$
    \item Update range: $\text{range} \gets \text{range} \ll 8$
\end{enumerate}

We have already consumed the first 4 bytes ($\text{0xA6}, \text{0x59}, \text{0x1C}, \text{0x00}$) as the initial code. The next byte is the fifth byte: $\text{0x00}$.

\textbf{Iteration 1:}
\[
\begin{aligned}
\text{low} &\gets (2{,}786{,}574{,}779 \ll 8) \mathrel{\&} 0\text{xFFFFFFFF} = 0\text{x1A2FBB00} = 439{,}463{,}680 \\
\text{code} &\gets ((2{,}789{,}021{,}696 \ll 8) \mathrel{\&} 0\text{xFFFFFFFF}) \mid 0 = 0\text{x591C0000} = 1{,}494{,}994{,}688 \\
\text{range} &\gets 5{,}153{,}960 \times 256 = 1{,}319{,}413{,}760
\end{aligned}
\]

Now $\text{range} = 1{,}319{,}413{,}760 \ge 16{,}777{,}216$, so stop renormalization.

\paragraph{Step 6: Decode Sixth Symbol}

Current interval: $[439{,}463{,}680, 1{,}758{,}877{,}440)$

Compute offset:
\[
u = \frac{\text{code} - \text{low}}{\text{range}} = \frac{1{,}494{,}994{,}688 - 439{,}463{,}680}{1{,}319{,}413{,}760} = \frac{1{,}055{,}531{,}008}{1{,}319{,}413{,}760} = 0.8
\]

Find symbol containing $0.8$:
\begin{itemize}
    \item $A$: $[0, 0.5)$ — no
    \item $B$: $[0.5, 0.8)$ — note: $B$ is $[0.5, 0.8)$, which \textbf{excludes} $0.8$
    \item $C$: $[0.8, 1.0)$ — contains $0.8$? Yes (the interval includes its left endpoint)
\end{itemize}
Sixth symbol is $C$.

\paragraph{Decoding Complete}

We have recovered the original sequence:
\[
\boxed{\texttt{B A C C C C}}
\]

\paragraph{Decoding State Table}

\begin{center}
\begin{tabular}{c|c|c|c|c}
\hline
\textbf{Step} & \textbf{Code} & \textbf{Low} & \textbf{Range} & \textbf{Symbol} \\ \hline
Start & $2{,}789{,}021{,}696$ & $0$ & $4{,}294{,}967{,}296$ & — \\
1 & $2{,}789{,}021{,}696$ & $2{,}147{,}483{,}648$ & $1{,}288{,}490{,}188$ & B \\
2 & $2{,}789{,}021{,}696$ & $2{,}147{,}483{,}648$ & $644{,}245{,}094$ & A \\
3 & $2{,}789{,}021{,}696$ & $2{,}662{,}879{,}723$ & $128{,}849{,}018$ & C \\
4 & $2{,}789{,}021{,}696$ & $2{,}765{,}958{,}937$ & $25{,}769{,}803$ & C \\
5 & $2{,}789{,}021{,}696$ & $2{,}786{,}574{,}779$ & $5{,}153{,}960$ & C \\
Renorm & $1{,}494{,}994{,}688$ & $439{,}463{,}680$ & $1{,}319{,}413{,}760$ & — \\
6 & $1{,}494{,}994{,}688$ & $439{,}463{,}680$ & $1{,}319{,}413{,}760$ & C \\ \hline
\end{tabular}
\end{center}

\paragraph{Key Observations}

\begin{enumerate}
    \item \textbf{Renormalization symmetry:} The encoder outputs a byte when $\text{range} < 2^{24}$; the decoder reads that same byte when $\text{range} < 2^{24}$ and incorporates it into $\text{code}$. Both update $\text{low}$ and $\text{range}$ identically.
    
    \item \textbf{Why this sequence triggers renormalization:} Four consecutive $C$'s (each with frequency $2/10 = 0.2$) cause the range to shrink by a factor of $0.2$ each time:
    \[
    644{,}245{,}094 \times 0.2 = 128{,}849{,}018
    \]
    \[
    128{,}849{,}018 \times 0.2 = 25{,}769{,}803
    \]
    \[
    25{,}769{,}803 \times 0.2 = 5{,}153{,}960 < 16{,}777{,}216
    \]
    
    \item \textbf{The mixed symbols:} Starting with $B$ and $A$ demonstrates different cumulative frequency calculations ($c_B=5$ for $B$, $c_A=0$ for $A$) before the repeated $C$'s trigger renormalization.
    
    \item \textbf{Interval visualization:} The interval narrowing table shows how each symbol selects a subinterval proportional to its frequency, progressively narrowing the range until renormalization resets it to a larger value.
    
    \item \textbf{Carry propagation:} In practice, implementations must handle carries when $\text{low}$ overflows during shifting. This is why real encoders use a buffer rather than outputting bytes immediately.
\end{enumerate}


\paragraph{Practical Implementation Note}

In actual range coding implementations, the final $\text{low}$ value is typically flushed by outputting the remaining bytes to ensure the decoder has enough information. This example shows the core mechanism: renormalization outputs bytes during encoding and consumes them during decoding when the range becomes too small to maintain precision.

\subsection{Asymmetric Numeral Systems (ANS): The Breakthrough}

Introduced in 2009 by Jarosław Duda, ANS revolutionized entropy coding by achieving:

\begin{itemize}
    \item compression efficiency very close to arithmetic coding,
    \item speeds approaching and often exceeding Huffman coding,
    \item simple integer arithmetic,
    \item elegant mathematical structure,
    \item ideal suitability for software and hardware implementations.
\end{itemize}

\subsubsection{Historical Context}

Before ANS, entropy coding was effectively split between:

\begin{itemize}
    \item \textbf{Huffman coding}: fast but suboptimal (integer bits only).
    \item \textbf{Arithmetic coding}: optimal but comparatively slow (requires multiplication/division).
\end{itemize}

Duda's insight was to invert the typical relationship between symbols and intervals, encoding information into the \emph{magnitude} of a single integer state.

\subsubsection{The Core Insight}

ANS is founded on a simple idea:

\begin{quote}
\emph{Maintain a single integer state $x$.  
Encoding a symbol transforms $x$ into a larger value $x'$ whose growth approximates the ideal information cost $-\log_2 p(s)$.}
\end{quote}

If a symbol $s$ has probability $p$, ANS ensures:

\[
\log_2(x') - \log_2(x) \approx -\log_2 p.
\]

This mirrors arithmetic coding but uses integer state transitions instead of intervals.

\subsubsection{Intuitive Explanation: The ``Stack'' Property}

Arithmetic coding emits bits from the most significant side.  
ANS, however, emits (and consumes) bits at the \emph{least significant side}, which reverses the order of symbols:

\begin{itemize}
    \item Encoding: symbols push bits onto the end of the state (like a stack).
    \item Decoding: symbols are read in reverse order, popping bits from the state.
\end{itemize}

This natural LIFO (Last-In-First-Out) behavior means:
\begin{itemize}
    \item If you encode $s_1, s_2, \ldots, s_n$ in that order, decoding yields $s_n, \ldots, s_2, s_1$
    \item To decode in natural order, either encode in reverse or buffer and reverse the output
\end{itemize}

\subsubsection{Why ANS Is Faster}

\begin{itemize}
    \item A single integer state variable (no interval maintenance).
    \item No divisions in the fast path for tANS (precomputed tables).
    \item Sequential memory access (no tree traversal like Huffman).
    \item Compatible with SIMD and GPU architectures.
\end{itemize}

ANS supports two major variants: rANS and tANS.

\subsection{Variations of ANS}

\subsubsection{rANS (Range ANS)}

rANS uses arithmetic-like interval partitioning:

\[
x' = \left\lfloor \frac{x}{f(s)} \right\rfloor \cdot \text{total} + (x \bmod f(s)) + F(s)
\]

where:
\begin{itemize}
    \item $x$ is the current state
    \item $f(s)$ is the frequency of symbol $s$
    \item $F(s)$ is the cumulative frequency of all symbols before $s$
    \item $\text{total} = \sum f(s)$ (typically a power of two)
\end{itemize}

Key properties:

\begin{itemize}
    \item Very simple to implement.
    \item Excellent compression efficiency.
    \item Naturally streaming (push/pop bits from LSB side).
\end{itemize}

\subsubsection{tANS (Table ANS)}

tANS precomputes a finite-state automaton (FSA) that maps:

\[
\text{state},\; \text{symbol} \longrightarrow \text{new state} + \text{emitted bits}.
\]

Properties:

\begin{itemize}
    \item Extremely fast (used in Zstandard).
    \item No multiplications or divisions at runtime.
    \item Ideal for hardware (lookup tables in ROM).
    \item Memory: approximately $4 \times \text{states}$ bytes (e.g., 16 KB for 4096 states).
\end{itemize}

\subsubsection{Comparison}

\begin{itemize}
    \item \textbf{rANS}: easier to implement, flexible, great for streaming, $O(1)$ memory.
    \item \textbf{tANS}: faster, more complex, uses $O(\text{states})$ memory, ideal for small alphabets.
\end{itemize}

\subsection{Detailed rANS Example}

To make the mechanics of rANS fully concrete, this section walks through a complete encoding and decoding example.  
This includes:

\begin{itemize}
    \item constructing normalized frequencies,
    \item building cumulative tables,
    \item understanding renormalization,
    \item encoding a symbol sequence step-by-step,
    \item decoding it back to the original sequence.
\end{itemize}

For pedagogical clarity, we use small numbers and a small symbol alphabet.  
Practical implementations use much larger states and ranges.

\subsubsection{Normalized Frequencies}

Given an alphabet $\{A, B\}$ with raw frequencies:

\[
f(A)=2,\qquad f(B)=1,\qquad \text{total}=3,
\]

we normalize the frequencies to a power-of-two total:

\[
L = 4.
\]

Normalized frequencies:

\[
f_A = 3,\qquad f_B = 1.
\]

Cumulative frequency ranges:

\[
A:\; [0,3),\qquad B:\; [3,4).
\]

Throughout this example:

\[
L = 4.
\]

\subsubsection{Initial State}

rANS maintains a single integer state $x$.  
To keep the example small, we set:

\[
x_0 = L = 4.
\]

Practical implementations initialize $x$ to a much larger value (e.g., $2^{23}$ for rANS32 or $2^{56}$ for rANS64), but using $x_0=L$ is perfectly valid for illustrating the algorithm.

\subsubsection{Encoding Equations}

Given the current state $x$ and symbol $s$ with frequency $f(s)$ and cumulative range $F(s)$, rANS encodes using:

\[
x' = \left\lfloor \frac{x}{f(s)} \right\rfloor \cdot L + (x \bmod f(s)) + F(s).
\]

This increases the magnitude of $x$ so that its growth approximates $-\log_2 p(s)$ bits of information.

\subsubsection{Renormalization}

Since $x$ grows during encoding, we must periodically emit lower bits to keep it within a machine word.

Define a renormalization lower bound (for educational purposes):

\[
\text{threshold} = L \cdot 256.
\]

While $x \geq \text{threshold}$:

\begin{enumerate}
    \item Emit the lowest 8 bits of $x$ (LSB-first).
    \item Set $x \leftarrow x \gg 8$.
\end{enumerate}

This keeps $x$ at a manageable size. (Production implementations use much larger thresholds.)

\subsubsection{Encoding the Sequence \texttt{B A B A}}

\textbf{Important}: Because ANS is a stack-like coder, we encode symbols in \textbf{reverse order} so that decoding produces the original order. For sequence \texttt{B A B A} (original order), we encode:

\[
\texttt{A},\; \texttt{B},\; \texttt{A},\; \texttt{B}
\]

\textbf{Step 1: Encode A}

Current state:

\[
x = 4.
\]

For $A$:

\[
f_A = 3,\qquad F(A)=0.
\]

Compute:

\[
x' = \left\lfloor \frac{4}{3} \right\rfloor \cdot 4 + (4 \bmod 3) + 0
    = 1 \cdot 4 + 1 = 5.
\]

No renormalization required.

\[
x = 5.
\]

\textbf{Step 2: Encode B}

\[
x = 5.
\]

For $B$:

\[
f_B = 1,\qquad F(B)=3.
\]

Compute:

\[
x' = \left\lfloor \frac{5}{1} \right\rfloor \cdot 4 + (5 \bmod 1) + 3
    = 5 \cdot 4 + 0 + 3 = 23.
\]

Still below the threshold:

\[
x = 23.
\]

\textbf{Step 3: Encode A}

\[
x = 23.
\]

For $A$:

\[
x' = \left\lfloor \frac{23}{3} \right\rfloor \cdot 4 + (23 \bmod 3)
    = 7 \cdot 4 + 2 = 30.
\]

\[
x = 30.
\]

\textbf{Step 4: Encode B}

\[
x = 30.
\]

For $B$:

\[
x' = \left\lfloor \frac{30}{1} \right\rfloor \cdot 4 + 3 = 30 \cdot 4 + 3 = 123.
\]

\[
x = 123.
\]

No renormalization was triggered because our threshold $=4\cdot256=1024$.

\subsubsection{Final Compressed State}

At the end of encoding:

\[
x_{\text{final}} = 123.
\]

The output stream consists of:

\begin{itemize}
    \item All emitted renormalization bytes (none in this example).
    \item The final state $x_{\text{final}}$ (which will be read by the decoder).
\end{itemize}

\subsubsection{Decoding}

Decoding processes symbols in the \textbf{forward} order (starting from the final state) and will produce symbols in reverse order of encoding. Since we encoded \texttt{A, B, A, B}, decoding will produce:

\[
\texttt{B},\; \texttt{A},\; \texttt{B},\; \texttt{A}
\]

which matches our original sequence \texttt{B A B A}.

Given current state $x$, we determine which symbol was encoded by computing:

\[
\text{slot} = x \bmod L.
\]

Then:

\[
\text{symbol} = \text{lookup\_by\_cumulative(slot)}.
\]

For our cumulative table:

\[
0,1,2 \rightarrow A,\qquad 3 \rightarrow B.
\]

After identifying the symbol $s$, we compute the previous state:

\[
x_{\text{prev}} = f(s) \cdot \left\lfloor \frac{x}{L} \right\rfloor + (x \bmod L) - F(s).
\]

Finally, if $x_{\text{prev}} < L$, we pull bytes from the input stream until:

\[
x_{\text{prev}} \geq L.
\]

\subsubsection{Decoding Step-by-Step}

\paragraph{Start}

\[
x = 123.
\]

\paragraph{Step 1}

\[
\text{slot} = 123 \bmod 4 = 3.
\]

Slot 3 corresponds to $B$.

Compute:

\[
x_{\text{prev}} = 1 \cdot \left\lfloor \frac{123}{4} \right\rfloor + 3 - 3
= 1 \cdot 30 + 0 = 30.
\]

\[
x = 30.
\]

\paragraph{Step 2}

\[
\text{slot} = 30 \bmod 4 = 2.
\]

Slot 2 → $A$.

\[
x_{\text{prev}} = 3 \cdot \left\lfloor \frac{30}{4} \right\rfloor + (2 - 0) = 3 \cdot 7 + 2 = 23.
\]

\[
x = 23.
\]

\paragraph{Step 3}

\[
\text{slot} = 23 \bmod 4 = 3.
\]

Slot 3 → $B$.

\[
x_{\text{prev}} = 1 \cdot \left\lfloor \frac{23}{4} \right\rfloor + (3 - 3) = 1 \cdot 5 + 0 = 5.
\]

\[
x = 5.
\]

\paragraph{Step 4}

\[
\text{slot} = 5 \bmod 4 = 1.
\]

Slot 1 → $A$.

\[
x_{\text{prev}} = 3 \cdot \left\lfloor \frac{5}{4} \right\rfloor + (1 - 0) = 3 \cdot 1 + 1 = 4.
\]

\[
x = 4.
\]

This matches $x_0 = 4$, confirming correct decoding.

\subsubsection{Summary}

The decoded symbol sequence (in order of decoding):

\[
\boxed{\texttt{B A B A}}
\]

matches the original.  
This demonstrates that:
\begin{itemize}
    \item Encoding in reverse order: \texttt{A, B, A, B}
    \item Decoding in forward order: \texttt{B, A, B, A}
    \item Result matches original: \texttt{B A B A}
\end{itemize}

\subsection{ANS in Modern Compressors}

ANS is widely used in state-of-the-art compression formats and libraries due to its exceptional performance and flexibility.

\subsubsection{Zstandard (Facebook / Meta)}

Zstandard's entropy stage uses tANS extensively.  
Benefits include:

\begin{itemize}
    \item extremely fast decoding,
    \item efficient table-driven implementation,
    \item well-suited for CPU vectorization,
    \item significantly better speed than arithmetic coding in practice.
\end{itemize}

tANS contributes significantly to Zstd's reputation as one of the fastest modern general-purpose compressors.

\subsubsection{Apple LZFSE}

LZFSE uses ANS as a core component to replace zlib-style Huffman coding.  
Advantages include:

\begin{itemize}
    \item improved compression ratio over zlib,
    \item much higher throughput,
    \item optimized for modern ARM processors.
\end{itemize}

\subsubsection{Google Draco}

Draco uses ANS for compressing 3D geometry and mesh connectivity.  
Reasons include:

\begin{itemize}
    \item compact symbol representation,
    \item suitability for GPU-friendly decoding,
    \item excellent branching behavior.
\end{itemize}

\subsubsection{JPEG XL}

JPEG XL employs ANS as part of its modular and VarDCT encoding pipelines.  
ANS allows:

\begin{itemize}
    \item high-speed entropy coding for transforms,
    \item efficient modeling of image coefficients,
    \item extremely competitive performance compared to arithmetic coding while avoiding multipliers.
\end{itemize}

\subsection{Hands-On: Implementing a Simple ANS Coder}

This section provides a complete pedagogical implementation of rANS in Python.  
Real-world implementations use optimizations, but this version is correct and readable.

\subsubsection{Educational rANS Implementation}

\begin{lstlisting}[language=Python, caption=Educational rANS Implementation]
class EducationalRANS:
    """Educational rANS implementation - demonstrates core concepts"""
    
    def __init__(self, freqs, L=1024):
        """
        Initialize rANS coder with symbol frequencies.
        
        Args:
            freqs: List of frequencies for symbols 0..n-1
            L: Total frequency (power of two, >= sum(freqs))
        """
        self.L = L
        # Normalize frequencies to sum to L
        total_raw = sum(freqs)
        self.freqs = []
        for f in freqs:
            # Ensure at least 1 occurrence
            norm = max(1, int(f * L / total_raw))
            self.freqs.append(norm)
        
        # Adjust to exactly sum to L (add residuals to largest symbols)
        diff = L - sum(self.freqs)
        for i in range(diff):
            self.freqs[i % len(freqs)] += 1
        
        # Build cumulative frequencies
        self.cum = [0]
        for f in self.freqs[:-1]:
            self.cum.append(self.cum[-1] + f)
        
        # Build lookup table for decoding
        self.slot_to_sym = [-1] * L
        for sym, (start, end) in enumerate(zip(self.cum, 
                                                [c + f for c, f in zip(self.cum, self.freqs)])):
            for slot in range(start, end):
                self.slot_to_sym[slot] = sym
    
    def encode(self, data):
        """
        Encode data to byte stream.
        Note: Encodes in reverse order due to ANS LIFO property.
        """
        x = self.L  # Initial state
        output = []
        
        # Encode in reverse order so decoding gives original order
        for sym in reversed(data):
            # Renormalization: emit bytes when state is large
            while x >= self.L * 1024:  # Example threshold
                output.append(x & 0xFF)
                x >>= 8
            
            # rANS encoding step
            freq = self.freqs[sym]
            cum = self.cum[sym]
            q = x // freq
            r = x % freq
            x = q * self.L + cum + r
        
        # Output final state
        while x > 0:
            output.append(x & 0xFF)
            x >>= 8
        
        return output
    
    def decode(self, byte_stream):
        """
        Decode from byte stream.
        Returns symbols in original order.
        """
        # Initialize state from bytes (LSB-first)
        x = 0
        for b in reversed(byte_stream):
            x = (x << 8) | b
        
        result = []
        
        while x > self.L:
            # rANS decoding step
            slot = x % self.L
            sym = self.slot_to_sym[slot]
            
            freq = self.freqs[sym]
            cum = self.cum[sym]
            x = freq * (x // self.L) + (slot - cum)
            
            result.append(sym)
        
        # Result is already in original order because we encoded in reverse
        return result

# Example usage
def test_rans():
    # Alphabet: A=0, B=1 with frequencies 2 and 1
    freqs = [2, 1]
    data = [1, 0, 1, 0]  # B A B A
    
    encoder = EducationalRANS(freqs, L=4)
    encoded = encoder.encode(data)
    decoded = encoder.decode(encoded)
    
    print(f"Original: {data}")
    print(f"Encoded bytes: {encoded}")
    print(f"Decoded: {decoded}")
    assert decoded == data
    print("Roundtrip successful!")

if __name__ == "__main__":
    test_rans()
\end{lstlisting}

\subsubsection{Key Implementation Notes}

\begin{itemize}
    \item \textbf{Reverse encoding}: The encode function processes symbols in reverse order so decode returns original order.
    \item \textbf{Normalization}: Frequencies are normalized to sum to a power of two for efficient division (using bit shifts in optimized implementations).
    \item \textbf{Renormalization threshold}: The example uses $L \times 1024$; production code uses much larger thresholds.
    \item \textbf{Slot lookup}: The decoder uses a precomputed array for $O(1)$ symbol lookup.
\end{itemize}

\subsection{Summary}

This chapter presented two powerful entropy coders that complement or replace traditional arithmetic coding:

\begin{itemize}
    \item \textbf{Range Coding} — an integer-based variant of arithmetic coding, widely used in LZMA, WebP, and ffmpeg.
    \item \textbf{ANS (Asymmetric Numeral Systems)} — a revolutionary coding system offering arithmetic-like compression efficiency with Huffman-like speed.
\end{itemize}

Two variants of ANS were covered:

\begin{itemize}
    \item \textbf{rANS} — streaming, simple, and ideal for software implementations.
    \item \textbf{tANS} — table-driven, extremely fast, ideal for high-performance applications.
\end{itemize}

Key takeaways:
\begin{enumerate}
    \item Range coding and ANS both achieve fractional bits per symbol (unlike Huffman).
    \item ANS's LIFO property is fundamental — information is pushed to the least significant side.
    \item rANS uses $O(1)$ memory; tANS uses $O(\text{states})$ memory but offers higher speed.
    \item Both are patent-free and widely used in modern compressors.
\end{enumerate}

\subsection{Exercises}

\begin{enumerate}
    \item \textbf{Range Coding Implementation:} Construct a range coder for an alphabet of size 3 using 32-bit arithmetic. Encode and decode the sequence \texttt{CABAC}. Show all intermediate steps including renormalization.
    
    \item \textbf{rANS Bits Emitted:} Derive the exact number of bits emitted by the rANS encoding of \texttt{B A B A} when using a practical renormalization threshold of $L \times 1024$. How does this compare to the theoretical entropy?
    
    \item \textbf{tANS Table Construction:} Build a tANS table for frequencies $\{A:5, B:3, C:2\}$ with $L=16$ states. Trace encoding and decoding of the sequence \texttt{ABCAB}.
    
    \item \textbf{Performance Comparison:} Implement rANS32 in C or Python and benchmark it against a standard Huffman coder on:
    \begin{itemize}
        \item Random data with uniform distribution
        \item English text (e.g., a book chapter)
        \item Highly skewed data ($p(A)=0.95$, others small)
    \end{itemize}
    Compare compression ratio and speed.
    
    \item \textbf{Mathematical Analysis:} Compare arithmetic coding and rANS for a skewed distribution (e.g., $p(A)=0.95$, $p(B)=0.05$). Which produces more renormalization steps? Explain why.
    
    \item \textbf{Adaptive ANS:} Extend the rANS implementation to support adaptive probabilities. Use a simple update rule (e.g., increment frequency after each symbol). Compare compression ratio against static ANS on non-stationary data.
    
    \item \textbf{LIFO Analysis:} Explain why ANS's LIFO property is problematic for streaming applications. Propose two different solutions and discuss their trade-offs.
\end{enumerate}

\subsection*{Additional Resources}

\begin{itemize}
    \item Duda, J. (2013). "Asymmetric numeral systems: entropy coding combining speed of Huffman coding with compression rate of arithmetic coding." arXiv:1311.2540.
    \item Giesen, F. (2014). "rANS in practice." \texttt{https://fgiesen.wordpress.com/}
    \item Collet, Y. (2016). "Finite State Entropy: A new breed of entropy coder." \texttt{https://github.com/Cyan4973/FiniteStateEntropy}
    \item Zstandard source code: \texttt{lib/compress/zstd\_opt.c} and \texttt{lib/entropy\_common.c}
    \item WebP source code: \texttt{https://chromium.googlesource.com/webm/libwebp/}
\end{itemize}
